\documentclass[11pt,a4paper]{article}

% ── Leichtes Preamble: kompiliert schnell und lokal (MiKTeX/pdflatex) ──
\usepackage[utf8]{inputenc}
\usepackage[T1]{fontenc}
\usepackage[ngerman]{babel}
\usepackage{lmodern}
\usepackage[a4paper,margin=2.5cm]{geometry}
\usepackage{enumitem}
\usepackage{booktabs}
\usepackage{tabularx}
\usepackage{array}
\usepackage{graphicx}
\usepackage{parskip}
\usepackage{url}

\newcolumntype{Y}{>{\raggedright\arraybackslash}X}
\setlist[itemize]{leftmargin=1.4em,itemsep=2pt,topsep=3pt}
\setlist[enumerate]{leftmargin=1.6em,itemsep=2pt,topsep=3pt}

% Platzhalter fuer Screenshots: das Dokument kompiliert OHNE Bilddateien.
% Spaeter einfach durch \includegraphics{...} ersetzen.
\newcommand{\platzhalter}[1]{%
  \begin{center}%
  \setlength{\fboxsep}{10pt}%
  \fbox{\parbox[c][3.2cm][c]{0.86\linewidth}{\centering\itshape #1}}%
  \end{center}%
}

% ══════════════════════════════════════════════════════════════════════════════
\begin{document}

% ── Titelseite ────────────────────────────────────────────────────────────────
\begin{titlepage}
\centering
\vspace*{2.2cm}
{\Large\scshape Technische Hochschule Mittelhessen}\\[2pt]
{\normalsize Fachbereich Mathematik, Naturwissenschaften und Informatik (MNI)}\\[0.4cm]
{\normalsize Praktikum Wirtschaftsinformatik --- Sommersemester 2026}\\[3.0cm]

\rule{\linewidth}{1pt}\\[8pt]
{\Huge\bfseries Aufgabendatenbank}\\[10pt]
{\LARGE Projektdokumentation}\\[8pt]
\rule{\linewidth}{1pt}\\[1.2cm]

{\large Ein Werkzeug zum Beschreiben, Strukturieren und Speichern\\von Lernaufgaben}\\[2.6cm]

\begin{tabular}{rl}
\textbf{Betreuung:} & Prof.\ Dr.\ Markus Siepermann\\
 & Johannes Kunz\\[6pt]
\textbf{Projektteam:} & [Vorname Nachname] \textit{(Projektleitung)}\\
 & [Vorname Nachname]\\
 & [Vorname Nachname]\\[6pt]
\textbf{Gruppe:} & [Gruppennummer]\\[6pt]
\textbf{Abgabe:} & 14.\ Juli 2026\\
\end{tabular}

\vfill
{\footnotesize Quelltext: GitHub-Repository \texttt{PWI-Aufgabendatenbank}}
\end{titlepage}

\tableofcontents
\newpage

% ══════════════════════════════════════════════════════════════════════════════
\section{Kurze Zusammenfassung}

Ausgangspunkt des Projekts war ein vorgegebenes relationales Datenbankschema für
eine \emph{Aufgabendatenbank}. Ziel war es, auf dieser Grundlage eine
vollständige, lauffähige Webanwendung zu entwickeln, mit der Lehrende Aufgaben
strukturiert erfassen, mit Metadaten versehen und untereinander in Beziehung
setzen können.

Umgesetzt wurde eine Drei-Schichten-Anwendung aus einem Vue-Frontend, einem
Spring-Boot-Backend und einer PostgreSQL-Datenbank. Alle zentralen fachlichen
Konzepte des Schemas sind implementiert: Aufgaben mit Metadaten (Autor, Lizenz,
Typ), Inhaltsbausteine in mehreren Formaten (Text, Bild, PDF), geordnete und
ungeordnete Kollektionen, horizontale Varianten, Validatoren, hierarchische Tags
mit Filter- und Suchfunktion sowie Darstellungs-Templates. Für den Betrieb auf
der E-Learning-Plattform der Hochschule wurde eine automatisierte Build-Pipeline
(CI\,$\rightarrow$\,GHCR) samt Konfigurationsdateien vorbereitet.

Das Projektteam bestand zu Beginn aus fünf Personen und wurde nach dem Ausstieg
zweier Mitglieder in der Anfangsphase auf drei Personen verkleinert; der
Funktionsumfang wurde vollständig auf dieses Team umverteilt und umgesetzt. Nicht
realisiert wurden das eigentliche Cluster-Deployment (das zuständige
Infrastruktur-Team stellte seine Arbeit ein) sowie eine reale Authentifizierung
und eine ursprünglich angedachte KI-Funktion (siehe Kapitel~\ref{sec:grenzen}).

% ══════════════════════════════════════════════════════════════════════════════
\section{Aufgabenstellung und Zielsetzung}

\subsection{Aufgabenstellung}
Die Aufgabe bestand darin, aus einem vorgegebenen relationalen Datenmodell eine
funktionsfähige Anwendung zu entwickeln. Das Schema definierte die zentralen
Entitäten und ihre Beziehungen; die konkrete Anwendung --- Frontend, Backend und
Persistenzschicht --- war vom Team zu entwerfen und zu implementieren. Es
handelte sich also nicht um ein Projekt „auf der grünen Wiese``, sondern um die
fachlich korrekte Umsetzung einer vorgegebenen Datenstruktur in eine benutzbare
Software.

\subsection{Abgrenzung des Funktionsumfangs}
Eine frühe, für das gesamte Projekt prägende Klärung mit dem Betreuer betraf die
Abgrenzung: Die Anwendung ist ausdrücklich ein \textbf{Autoren- und
Verwaltungswerkzeug} und \emph{kein} Lern- oder Prüfungssystem. Sie unterstützt
das Erstellen und Organisieren von Aufgaben, nicht deren Bearbeitung durch
Lernende.

\begin{itemize}
  \item \textbf{Im Verantwortungsbereich:} Aufgaben erfassen, mit Metadaten
  versehen, in Sequenzen und Varianten in Beziehung setzen sowie Prüfregeln
  (Validatoren) \emph{definieren und speichern}.
  \item \textbf{Nicht im Verantwortungsbereich:} das eigentliche Lösen oder
  Korrigieren von Aufgaben sowie die automatische Erzeugung von Varianten. Diese
  Ausführung übernimmt ein separater Dienst der Plattform.
  \item \textbf{Zugriffsmodell:} In dieser Iteration gibt es bewusst keine
  Rechteeinschränkung; jeder Nutzer kann alles anlegen und bearbeiten. Eine
  spätere zentrale Anmeldung (THM-CAS) kann darauf aufsetzen.
\end{itemize}

\subsection{Zielsetzung}
Angestrebt wurde eine fachlich korrekte, technisch saubere und erweiterbare
Umsetzung des Datenmodells mit einer benutzerfreundlichen Oberfläche: alle
Kern-Metadatenkonzepte abbilden, die beiden im Schema angelegten Beziehungsarten
(vertikal und horizontal) unterstützen und die Anwendung für den Betrieb auf der
Hochschulplattform vorbereiten.

% ══════════════════════════════════════════════════════════════════════════════
\section{Ausgangssituation und Recherche}

\subsection{Ausgangssituation}
Zu Projektbeginn lagen das relationale Schema und eine grobe fachliche
Beschreibung vor. Die erste Aufgabe war daher weniger das Programmieren als das
\emph{Verstehen} des Modells. Zwei Punkte erwiesen sich als zentral und mussten
im Team gemeinsam erarbeitet werden:

\begin{itemize}
  \item \textbf{Eine Kollektion ist selbst ein Item.} Es gibt keine getrennte
  „Sammlung`` neben der „Aufgabe``; über \texttt{parent\_item\_id} \emph{ist}
  eine Kollektion zugleich ein Item und teilt sich dessen Metadaten und
  Verknüpfungen.
  \item \textbf{Es gibt zwei getrennte Beziehungsarten.} Das Schema kennt eine
  \emph{vertikale} (geordnete Sequenzen) und eine \emph{horizontale} (Varianten
  über ein gemeinsames Wurzel-Item) Beziehung. Beide dürfen nicht vermischt
  werden --- eine Erkenntnis, die später eine wichtige Entwurfsentscheidung nach
  sich zog (Abschnitt~\ref{sec:beziehungen}).
\end{itemize}

\subsection{Technologieauswahl und Begründung}\label{sec:tech}
Auf Basis dieses Verständnisses wurden geeignete Technologien recherchiert. Die
Auswahl orientierte sich an drei Kriterien: Eignung für die fachliche Aufgabe,
durchgängige Typsicherheit sowie gute Container- und Cluster-Tauglichkeit, da die
Plattform der Hochschule auf Kubernetes aufbaut. Neben der reinen Eignung waren
zwei praktische Gründe ausschlaggebend: der Betreuer stellte eine
\textbf{Vue-Vorlage} als Startpunkt bereit, und das Team hatte bereits
\textbf{Vorerfahrung} mit dem gewählten Stack (Vue, Spring Boot, PostgreSQL).

\begin{center}
\begin{tabularx}{\linewidth}{@{}l l Y@{}}
\toprule
\textbf{Technologie} & \textbf{Version} & \textbf{Eingesetzt für} \\
\midrule
Vue (Composition API) & 3.5.34 & reaktives Komponenten-Frontend \\
Vuetify & 4.0.7 & Material-Design-UI (Formulare, Dialoge) \\
Pinia & 3.0.4 & zentraler, typisierter Zustand (Store) \\
vuedraggable & 4.1.0 & Drag-and-Drop im Strukturbaum \\
CodeMirror 6 (vue-codemirror) & 6.x & XML-Editor für Darstellungs-Templates \\
axios & 1.6 & HTTP-Kommunikation mit dem Backend \\
Vite / Vitest & 8.0.12 / 4.1.6 & Build-Werkzeug / Test-Framework \\
TypeScript & 6.0 & Typsicherheit über alle Frontend-Schichten \\
Spring Boot & 3.2.5 & REST-Backend (Controller/Service/Repository) \\
Java & 21 & Backend-Sprache \\
Spring Data JPA & (Boot 3.2.5) & deklarative Persistenz auf dem Schema \\
PostgreSQL & 16 & relationale Datenbank \\
Docker Compose & --- & reproduzierbarer lokaler Start \\
\bottomrule
\end{tabularx}
\end{center}

\textbf{Frontend --- Vue, Vuetify, Pinia.} Vue~3 mit der Composition API bietet
ein reaktives Komponentenmodell, das sich gut für die zentrale Baum-/Editor-
Ansicht eignet. Ausschlaggebend waren zusätzlich die vom Betreuer bereitgestellte
Vue-Vorlage und die vorhandene Team-Erfahrung. Vuetify liefert eine konsistente
Komponentenbibliothek und reduziert den Aufwand für Formulare und Dialoge; Pinia
hält den gesamten Anwendungszustand in einem typisierten Store. Für zwei
spezielle Anforderungen kamen dedizierte Bibliotheken zum Einsatz:
\textbf{vuedraggable} für die Drag-and-Drop-Umsortierung im Strukturbaum und
\textbf{CodeMirror~6} (über \texttt{vue-codemirror}) als Editor für die
XML-Darstellungs-Templates. Alternativen wie React oder Angular wurden verworfen:
React hätte mehr manuelle Zustandsarchitektur erfordert, Angular wäre für den
Projektumfang überdimensioniert gewesen --- und beide hätten weder die Vorlage
noch die Vorerfahrung genutzt.

\textbf{Backend --- Spring Boot, Java 21.} Die Wahl von Spring Boot war eine
bewusste, auf die konkrete Anwendung zugeschnittene Entscheidung: Das Schema ist
klar relational, und Spring Boot bietet mit der konventionsbasierten Struktur
(Controller\,$\rightarrow$\,Service\,$\rightarrow$\,Repository) und \textbf{Spring
Data JPA} eine deklarative Persistenzschicht, die die vorgegebenen Tabellen
direkt auf Entitäten abbildet. \textbf{Spring Validation} sichert Eingaben auf
Controller-Ebene ab. Auch hier war die Vorerfahrung des Teams ein Vorteil, sodass
sich die Arbeit auf die Fachlogik konzentrieren konnte statt auf Boilerplate.

\textbf{Datenbank --- PostgreSQL.} PostgreSQL war durch das Schema faktisch
vorgezeichnet und passte fachlich sehr gut: native \texttt{UUID}-Primärschlüssel,
der \texttt{JSONB}-Typ (mit GIN-Index) für flexible Inhaltsbausteine sowie
\texttt{BYTEA} für Binärdaten (Bilder, PDF). Auch mit PostgreSQL hatte das Team
bereits Erfahrung.

% ══════════════════════════════════════════════════════════════════════════════
\section{Anforderungen und Rahmenbedingungen}\label{sec:anf}

\subsection{Funktionale Anforderungen}
\begin{enumerate}[label=\textbf{F\arabic*}]
  \item Aufgaben über ein Formular mit bewusster Bestätigung erstellen.
  \item Inhaltsbausteine unterschiedlicher Formate (Text, Bild, PDF) hinzufügen
  und anzeigen.
  \item Aufgaben zu geordneten Sequenzen (Lernpfaden) und ungeordneten
  Sammlungen (Kollektionen) zusammenfassen.
  \item Horizontale Varianten einer Aufgabe über ein gemeinsames Wurzel-Item
  darstellen.
  \item Validatoren (Restriktionen) definieren, speichern und Aufgaben zuweisen.
  \item Referenzdaten (Autor, Lizenz, Typ, Inhaltstyp) auswählen und neu anlegen.
  \item Hierarchische Tags zuweisen, anzeigen sowie zur Filterung und Suche
  nutzen.
  \item Darstellungs-Templates (XML) je Aufgabe bearbeiten und in einer Vorschau
  prüfen.
\end{enumerate}

\subsection{Nicht-funktionale Anforderungen}
\begin{enumerate}[label=\textbf{N\arabic*}]
  \item \textbf{Datenintegrität:} referentielle Konsistenz auf Datenbankebene.
  \item \textbf{Typsicherheit:} durchgehend typisierte Schnittstellen zwischen
  Frontend und Backend.
  \item \textbf{Wartbarkeit:} klare Schichtung, austauschbare Datenzugriffsschicht.
  \item \textbf{Reaktionsfreudigkeit:} optimistische UI-Aktualisierung für einen
  flüssigen Arbeitsablauf.
  \item \textbf{Portabilität:} reproduzierbarer Start über Container, vorbereitet
  für den Cluster-Betrieb.
\end{enumerate}

\subsection{Rahmenbedingungen}
\begin{itemize}
  \item \textbf{Plattform:} Die Anwendung soll auf der E-Learning-Plattform der
  THM laufen. Diese integriert Dienste über Kubernetes; der Zugriff erfolgt über
  ein zentrales API-Gateway (Traefik), die Images liegen in der GitHub Container
  Registry (GHCR).
  \item \textbf{Authentifizierung:} nach Plattform-Vorgabe zunächst zu
  \emph{mocken}; eine globale Anmeldung (THM-CAS) sollte später folgen.
  \item \textbf{Vorgegebenes Datenmodell:} Das Schema war gesetzt und durfte in
  seiner Grundstruktur nicht verändert werden.
  \item \textbf{Team und Zusammenarbeit:} Entwicklung im Team mit geschütztem
  Hauptzweig; jede Änderung wird über Pull-Requests begutachtet.
\end{itemize}

% ══════════════════════════════════════════════════════════════════════════════
\section{Konzept und Entwurf}\label{sec:konzept}

\subsection{Systemarchitektur}
Die Anwendung folgt einer klassischen Drei-Schichten-Architektur mit klarer
Trennung von Präsentation, Geschäftslogik und Persistenz. Frontend (Port 8085)
und Backend (Port 8080) kommunizieren zustandslos über HTTP/JSON; das Backend
greift über JPA auf die PostgreSQL-Datenbank (Port 5432) zu.

\begin{verbatim}
Adb.vue (Wurzel)
 |- AdbToolbar.vue      Aktionen (Aufgabe/Kollektion erstellen)
 |- AdbStructure.vue    Strukturbaum + Filterleiste (oben/links)
 |   \- AdbTreeItem.vue  rekursive Baum-Knoten (Ordner = Kollektion,
 |                       Datei = Aufgabe)
 \- AdbEditor.vue       Detailansicht (rechts)
     |- AdbTagsSection.vue      Tags der Aufgabe
     |- AdbContentList.vue      Inhaltsbausteine (Text/Bild/PDF)
     |- AdbValidatorEditor.vue  Validatoren
     \- AdbVariantsPanel.vue    horizontale Varianten

exerciseStore.ts (Pinia)  ->  ApiAdapter (Interface)
                                |- AdbApiService     (axios, echte API)
                                \- DevAdbApiService  (lokale Mock-Daten)
\end{verbatim}

\textbf{Adapter-Muster als Entkopplung.} Ein zentrales Entwurfselement ist das
Adapter-Muster: Der Store spricht ausschließlich gegen die Schnittstelle
\texttt{ApiAdapter}. Zwei Implementierungen erfüllen diesen Vertrag ---
\texttt{AdbApiService} für echte HTTP-Aufrufe und \texttt{DevAdbApiService} mit
lokalen Mock-Daten. Dadurch lässt sich das Frontend vollständig ohne laufendes
Backend entwickeln und vorführen (Dummy-Modus). Die Auswahl erfolgt zur Startzeit
über eine Umgebungsvariable. Das Backend ist spiegelbildlich in drei Ebenen
gegliedert: \textbf{Controller} (REST-Endpunkte, Eingabevalidierung),
\textbf{Service} (Geschäftslogik, Transaktionsgrenzen) und \textbf{Repository}
(Spring Data JPA). Die Datenübertragung erfolgt über typisierte DTOs; eine
gemeinsame Typdatei im Frontend spiegelt die Vertragsstruktur, sodass Abweichungen
zwischen Client und Server bereits zur Übersetzungszeit auffallen.

\subsection{Datenmodell: die Entitäten}
Das Schema umfasst 18 Tabellen in drei logischen Ebenen. Jede wird im Frontend
über den Store und die Editor-Komponenten bedient und im Backend durch eine
JPA-Entität, ein Repository und (soweit fachlich nötig) einen Service abgebildet.

\subsubsection{Referenz-Entitäten (Nachschlagelisten)}
Diese Tabellen liefern die auswählbaren Stammdaten. Im Backend sind es einfache
Entitäten mit Repository; im Frontend werden sie als Auswahllisten geladen.
\begin{itemize}
  \item \texttt{author} --- Autor eines Items oder Inhalts (Anzeige-Deskriptor,
  optional Mail).
  \item \texttt{license} --- Lizenz (z.\,B.\ CC-BY), an Item und Inhalt hängbar.
  \item \texttt{item\_type} --- Typ einer Aufgabe (fachliche Kategorie).
  \item \texttt{item\_content\_type} --- Typ eines Inhaltsbausteins (Text, Bild,
  PDF\,\ldots).
  \item \texttt{item\_representation\_template} --- XML-Vorlage, die beschreibt,
  wie die Inhalte einer Aufgabe angeordnet dargestellt werden.
  \item \texttt{tag} --- hierarchisches Schlagwort (Selbstreferenz
  \texttt{parent\_tag\_id}).
  \item \texttt{validator} --- methodische Restriktion (z.\,B.\ „muss INNER JOIN
  enthalten``).
  \item \texttt{modifier} --- Regel zur (automatischen) Variantenerzeugung
  \emph{(modelliert, in dieser Iteration ohne UI --- siehe
  Kapitel~\ref{sec:grenzen})}.
\end{itemize}

\subsubsection{Kern-Entitäten}
\textbf{\texttt{item} (Aufgabe).} Jede Aufgabe ist ein \texttt{item} mit
Pflichtbezügen zu Autor, Lizenz und Typ. Bemerkenswert: es gibt \emph{kein}
eigenes Titelfeld --- der Titel lebt als Inhaltsbaustein und ist über die
Verknüpfung mit einem \texttt{purpose} (z.\,B.\ „Aufgabenstellung``) zugeordnet.
Die Selbstreferenz \texttt{root\_item\_id} verbindet zusammengehörige Varianten.
Im Frontend ist ein Item ein Knoten im Strukturbaum; im Backend kapselt
\texttt{ItemService} das Anlegen, Ändern, Löschen und die Umwandlung in eine
Kollektion.

\textbf{\texttt{item\_content} (Inhaltsbaustein).} Trägt entweder strukturierte
Daten (\texttt{JSONB}) oder Binärdaten (\texttt{BYTEA}, für Bild/PDF) und besitzt
eigene Lizenz, Typ und Autor. Im Frontend wird ein Inhalt im
\texttt{AdbContentEditor} bearbeitet; Binärdaten werden per Multipart-Upload
übertragen.

\textbf{\texttt{item\_collection} (Kollektion).} Gruppiert Items. Das boolesche
Feld \texttt{"order"} unterscheidet geordnete Sequenzen (Sub-Items mit
\texttt{position} 1, 2, 3\,\ldots) von ungeordneten Sammlungen (\texttt{position}
= \texttt{NULL}). Über \texttt{parent\_item\_id} \emph{ist} eine Kollektion
zugleich ein Item und besitzt keine eigene Metadaten-Kopie.

\subsubsection{Verknüpfungs-Entitäten (n:m-Beziehungen)}
\begin{itemize}
  \item \texttt{item\_contents} --- ordnet einem Item seine Inhaltsbausteine mit
  \texttt{purpose} zu (eigene Verknüpfungsentität mit zusammengesetztem Schlüssel).
  \item \texttt{item\_collection\_sub\_item} --- ordnet einer Kollektion ihre
  Kind-Items samt \texttt{position} zu (die vertikale Beziehung).
  \item \texttt{item\_tags} / \texttt{item\_content\_tags} --- Tags an Items bzw.\
  an Inhalte.
  \item \texttt{item\_validator} --- Validatoren an Items.
  \item \texttt{item\_modifier} --- Modifikatoren an Items \emph{(vorbereitet)}.
  \item \texttt{item\_content\_types} --- erlaubte Inhaltstypen je Item.
\end{itemize}

\subsection{Beziehungskonzept: vertikal und horizontal}\label{sec:beziehungen}
Das Schema kennt bewusst \emph{zwei orthogonale} Beziehungsarten, die im Projekt
klar getrennt behandelt werden:

\begin{center}
\begin{tabularx}{\linewidth}{@{}l l Y@{}}
\toprule
 & \textbf{Vertikal (Sequenz)} & \textbf{Horizontal (Variante)} \\
\midrule
Bedeutung & aufeinander aufbauende, \emph{verschiedene} Aufgaben in fester
Reihenfolge & \emph{dieselbe} Aufgabe in mehreren Ausprägungen \\
Mechanismus & geordnete \texttt{item\_collection} + \texttt{position} &
gemeinsames \texttt{root\_item\_id} \\
Rolle des Items & Behälter (Lernpfad) & einzelne Aufgabe \\
\bottomrule
\end{tabularx}
\end{center}

Diese Trennung ist eine tragende Entwurfsentscheidung: Ein Item spielt zu einem
Zeitpunkt genau \emph{eine} Rolle. Beim Umwandeln einer Aufgabe mit Varianten in
eine Kollektion warnt die Anwendung daher ausdrücklich, statt beide
Beziehungsarten stillschweigend zu vermischen. Konsequent dazu lässt das Löschen
einer Kollektion die enthaltenen Aufgaben unberührt (sie werden nur aus der
Kollektion gelöst) und tastet ihren Varianten-Bezug nicht an.

\subsection{Referentielle Integrität}
Die Fremdschlüssel setzen bewusst unterschiedliche Löschstrategien ein:

\begin{center}
\begin{tabularx}{\linewidth}{@{}l l Y@{}}
\toprule
\textbf{Beziehung} & \textbf{Strategie} & \textbf{Wirkung} \\
\midrule
\texttt{item} $\rightarrow$ author/license/type & \texttt{RESTRICT} &
Stammdaten sind gegen Löschen geschützt \\
\texttt{item.root\_item\_id} $\rightarrow$ item & \texttt{SET NULL} & Varianten
überleben das Löschen des Wurzel-Items \\
\texttt{tag.parent\_tag\_id} $\rightarrow$ tag & \texttt{SET NULL} & Unter-Tags
werden zu Wurzel-Tags \\
\texttt{item\_collection.parent\_item\_id} $\rightarrow$ item & \texttt{CASCADE}
& die Kollektion verschwindet mit ihrem Item \\
\texttt{item\_tags}, \texttt{item\_contents}, \ldots & \texttt{CASCADE} &
Verknüpfungen werden mit Item/Content aufgeräumt \\
\bottomrule
\end{tabularx}
\end{center}

% ══════════════════════════════════════════════════════════════════════════════
\section{Umsetzung}\label{sec:umsetzung}

Dieser Abschnitt beschreibt die umgesetzten Funktionen und --- wo es das
Verständnis fördert --- was bei einer Aktion zwischen Frontend, Backend und
Datenbank tatsächlich passiert.

\subsection{Aufgabe erstellen (durchgängiger Ablauf)}
Statt eine Aufgabe sofort beim Klick anzulegen, öffnet die Anwendung einen
Formular-Dialog (Typ, Autor, Lizenz, Aufgabentext). Erst mit „Erstellen`` startet
der Ablauf: Der \textbf{Store} legt die Aufgabe \emph{optimistisch} sofort lokal
im Baum an (die UI reagiert unmittelbar) und ruft über den \textbf{Adapter}
nacheinander \texttt{POST /items} (Item anlegen), \texttt{POST /items/\{id\}/}
\texttt{contents} (Aufgabentext als Inhalt) und --- falls in einer Kollektion ---
\texttt{POST /collections/\{id\}/items} (Zuordnung) auf. Im \textbf{Backend}
nehmen \texttt{ItemController} und \texttt{ItemContentController} die Anfragen
entgegen, die zugehörigen Services führen die Fachlogik aus und die Repositories
schreiben in PostgreSQL. Die Antwort liefert die echten IDs, die der Store gegen
die temporären optimistischen IDs austauscht. Schlägt ein Schritt fehl, gleicht
der Store den betroffenen Teilbaum wieder mit dem Datenbankstand ab, sodass keine
„Geister-Aufgabe`` zurückbleibt.

\subsection{Inhalte in mehreren Formaten}
Inhaltsbausteine unterstützen Text (JSON), Bilder und PDF. Binärdaten werden per
Multipart-Upload übertragen und im Feld \texttt{blob\_serialized\_content}
abgelegt; Bilder erscheinen als Vorschau, PDFs als Download-Link. Jeder Baustein
trägt einen \texttt{purpose} und eigene Metadaten (Typ, Lizenz).

\subsection{Kollektionen und Reihenfolge}
Aufgaben lassen sich in Kollektionen gruppieren. Ein Umschalter überführt eine
Kollektion zwischen geordnet und ungeordnet; beim Einschalten vergibt das Backend
fortlaufende Positionen, beim Ausschalten werden sie auf \texttt{NULL} gesetzt
(Frontend und Backend werden dabei konsistent gehalten). Die Umsortierung
innerhalb einer geordneten Kollektion erfolgt per Drag-and-Drop
(\texttt{vuedraggable}); die Positionen werden serverseitig neu berechnet.

\subsection{Horizontale Varianten}
Varianten einer Aufgabe teilen sich ein \texttt{root\_item\_id} und werden im
\texttt{AdbVariantsPanel} angezeigt. Beim Auswählen einer Aufgabe lädt die
Anwendung automatisch alle Geschwister-Varianten des gemeinsamen Wurzel-Items.
Der oben genannte Warnhinweis verhindert, dass eine Aufgabe versehentlich
gleichzeitig Variante und Kollektion wird.

\subsection{Validatoren}
Validatoren beschreiben methodische Restriktionen. Sie werden gemäß der
Projektabgrenzung nur \emph{definiert, gespeichert und Aufgaben zugewiesen}; die
Ausführung erfolgt extern. Umgesetzt sind Anlegen, Bearbeiten, Löschen und
Zuordnung (\texttt{AdbValidatorEditor} $\rightarrow$ \texttt{ValidatorController}).

\subsection{Referenzdaten auswählen und anlegen}
Autor, Lizenz, Typ und Inhaltstyp werden aus dem Backend geladen und im Editor
über Auswahllisten gesetzt. Über eine „+~Neu``-Funktion lassen sie sich direkt
neu anlegen; doppelte Namen werden server- und clientseitig abgefangen.

\subsection{Hierarchische Tags, Filter und Suche}
Tags sind über \texttt{parent\_tag\_id} hierarchisch und werden in
Pfadschreibweise dargestellt (z.\,B.\ \texttt{Mathematik/Analysis/Ableitungen}).
Die Zuweisung erfolgt über einen aufklappbaren Baum mit Mehrfachauswahl; feste
Regeln gelten beim Anlegen (global eindeutig, keine Leerzeichen, genau ein
Elternknoten, kein gleichzeitiges Markieren von Tag und Vorfahr/Nachfahr). Eine
gemeinsame obere Filterleiste kombiniert den Tag-Filter (\emph{mit Vererbung}:
ein Eltern-Tag findet auch Aufgaben mit Nachfahr-Tags) mit der Text- und
Attributsuche (Typ, Autor).

\subsection{Darstellungs-Templates}
Je Aufgabe kann ein XML-Template die Anordnung der Inhalte beschreiben. Es wird im
\textbf{CodeMirror}-Editor bearbeitet, gegen die vorhandenen \texttt{purpose}s
validiert (fehlende/überzählige Bausteine werden gemeldet) und in einer Vorschau
geprüft.

\subsection{Qualitätssicherung}
Die Qualitätssicherung ruht auf drei Säulen:
\begin{itemize}
  \item \textbf{Statische Typprüfung:} \texttt{vue-tsc} im Frontend und der
  Java-Compiler im Backend fangen Vertragsverletzungen früh ab.
  \item \textbf{Automatisierte Tests:} Unit-Tests mit \texttt{vitest} (Frontend,
  u.\,a.\ für Store-Logik und die neu abgesicherten Fehlerfälle) sowie
  Mockito-basierte Service- und Controller-Tests im Backend --- ohne Bedarf an
  einer laufenden Datenbank.
  \item \textbf{Strukturvalidierung:} eine reine Funktion prüft den geladenen
  Baum gegen fachliche Regeln (nur Kollektionen dürfen Kinder haben, jede
  \texttt{root\_item\_id} muss auflösbar sein).
\end{itemize}

\subsection{Grundlage für den Cluster-Betrieb}
Für den Betrieb auf der Hochschulplattform ist die Anwendung als klassischer,
synchroner Microservice ausgelegt. Eine CI/CD-Pipeline (GitHub Actions) baut bei
jedem erfolgreichen Merge in den Hauptzweig die Container-Images für Backend und
Frontend und veröffentlicht sie in der GitHub Container Registry (GHCR). Je
Anwendung beschreibt eine \texttt{values.yaml}, wie der Dienst im Cluster heißt,
wie viele Repliken er benötigt und über welchen Netzwerkpfad er erreichbar ist;
\textbf{ArgoCD} sollte diese Konfiguration mit dem Cluster synchronisieren (zum
tatsächlichen Deployment siehe Kapitel~\ref{sec:grenzen}).

% ══════════════════════════════════════════════════════════════════════════════
\section{Projektorganisation}\label{sec:orga}

\subsection{Vorgehen und Versionsverwaltung}
Die Entwicklung erfolgte git-basiert mit Feature-Zweigen und Pull-Requests. Der
Hauptzweig ist gegen direkte Pushes geschützt; jede Änderung wurde vor dem
Zusammenführen von einem anderen Teammitglied begutachtet. Die Arbeit wurde über
ein \textbf{GitHub-Project-Board} organisiert, auf dem Aufgaben als Tickets durch
die Spalten \emph{In Progress}, \emph{Reviewed} und \emph{Closed} wanderten ---
so waren Stand und Zuständigkeiten jederzeit sichtbar.

\platzhalter{[Screenshot: GitHub-Project-Board (Spalten In Progress / Reviewed /
Closed) hier einfügen]}

\subsection{Zusammenarbeit im Team}
Über die reine Werkzeugnutzung hinaus prägten regelmäßige Absprachen die
Zusammenarbeit:
\begin{itemize}
  \item \textbf{Meetings ein- bis zweimal pro Woche} --- je nach Fortschritt und
  offenen Fragen. Sie dienten der Abstimmung, der Klärung von Schnittstellen und
  der gemeinsamen Priorisierung.
  \item \textbf{Zwei Pair-Programming-Sitzungen} --- für besonders verzahnte oder
  kniffelige Aufgaben, bei denen gemeinsames Arbeiten Fehler früh vermied und
  Wissen im Team verteilte.
\end{itemize}
Wiederkehrende Herausforderung war das Zusammenführen länger laufender
Feature-Zweige, da mehrere Funktionen dieselben zentralen Dateien (Store, Adapter,
Typen) berührten. Als Konsequenz wurde ein Arbeitsablauf mit häufiger,
frühzeitiger Integration und konsequenter Typ- und Testprüfung nach jedem Merge
etabliert.

\subsection{Projektplanung: Phasen und Meilensteine}
Das Projekt lief von Ende April bis zur Abgabe am 14.\ Juli 2026 (Aktivitäts-
schwerpunkt im Juni).

\begin{center}
\begin{tabularx}{\linewidth}{@{}l l Y@{}}
\toprule
\textbf{Phase} & \textbf{Zeitraum} & \textbf{Inhalt} \\
\midrule
Konzeption \& Einarbeitung & Ende Apr.\ -- Mitte Mai & Team-Aufbau, Analyse des
Schemas, Technologieauswahl, Projektstruktur \\
Neuausrichtung & Mitte Mai & Ausstieg zweier Mitglieder, Umverteilung auf drei
Personen \\
Grundgerüst & Mai & Drei-Schichten-Aufbau, DB-Schema, Docker-Compose,
Adapter-Muster, Baum-/Editor-Ansicht \\
Kernfunktionen & Juni & Inhalte, Kollektionen/Reihenfolge, Varianten,
Validatoren, Referenzdaten, Frontend-Backend-Anbindung \\
Erweiterung \& Integration & Juli & hierarchische Tags + Filter, Suche,
Repräsentations-Templates, CI\,$\rightarrow$\,GHCR, Robustheits-Fixes \\
Abgabe & 14.\ Juli & Dokumentation, Tests, finale Bereinigung \\
\bottomrule
\end{tabularx}
\end{center}

Wesentliche Meilensteine: die abgestimmte fachliche Konzeption, die erste
durchgängige Frontend-Backend-Anbindung, die fachliche Vollständigkeit der
Kernfunktionen sowie die einsatzbereite Build-Pipeline mit veröffentlichten
Images in GHCR.

\subsection{Aufgabenverteilung}
Die Arbeit war zunächst nach Schwerpunkten aufgeteilt, verlief zum Projektende hin
aber zunehmend übergreifend:

\begin{center}
\begin{tabularx}{\linewidth}{@{}l Y@{}}
\toprule
\textbf{Person} & \textbf{Schwerpunkt} \\
\midrule
{}[Vorname Nachname] \textit{(Projektleitung)} & Datenbank und Backend zu großen
Teilen, Deployment (CI/CD, GHCR, \texttt{values.yaml}) eigenständig,
Koordination; gegen Ende zusätzlich Frontend \\
{}[Vorname Nachname] & Schwerpunkt Frontend (Baum, Editor, UI-Komponenten); gegen
Ende auch Backend-Aufgaben \\
{}[Vorname Nachname] & Schwerpunkt Backend (Endpunkte, Persistenz); gegen Ende
auch Frontend-Aufgaben \\
\bottomrule
\end{tabularx}
\end{center}

Anfangs bearbeitete eine Person überwiegend das Frontend und berührte das Backend
nicht, während eine andere überwiegend im Backend arbeitete; die Projektleitung
verantwortete Datenbank, Backend und das gesamte Deployment eigenständig. Im
weiteren Verlauf lösten sich diese Grenzen auf --- teils als Folge der
verkleinerten Teamgröße, was zugleich dazu führte, dass alle drei Personen ein
gutes Gesamtverständnis des Systems erlangten.

\subsection{Abweichungen vom ursprünglichen Plan}
Das Projekt startete mit \textbf{fünf} Personen. Nach etwa drei Wochen stiegen
zwei Mitglieder aus, sodass die verbleibende Arbeit auf \textbf{drei} Personen neu
verteilt werden musste. Das erhöhte die Last pro Person spürbar und war der
wesentliche Grund für die spätere schichtübergreifende Arbeitsweise. Weitere
inhaltliche Abweichungen sind in Kapitel~\ref{sec:grenzen} zusammengefasst.

% ══════════════════════════════════════════════════════════════════════════════
\section{Grenzen und Ausblick}\label{sec:grenzen}

Die folgende ehrliche Bestandsaufnahme trennt, was umgesetzt wurde, was bewusst
zurückgestellt und was aus externen Gründen nicht realisiert werden konnte.

\subsection{Nicht realisiert (mit Begründung)}
\begin{itemize}
  \item \textbf{Cluster-Deployment.} Die Build-Pipeline (CI\,$\rightarrow$\,GHCR)
  und die \texttt{values.yaml} je App sind fertig. Das eigentliche Ausrollen im
  Cluster kam jedoch nicht zustande: Auf die Anfrage, uns in ArgoCD einzutragen,
  teilte das zuständige Infrastruktur-Team (Gruppe~1) mit, dass es das Projekt
  \emph{nicht weiterführt}. Ohne diese Registrierung war ein Deployment nicht
  möglich --- ein externer, nicht vom Team zu verantwortender Umstand.
  \item \textbf{Authentifizierung.} Nach Plattform-Vorgabe sollte die
  Authentifizierung zunächst gemockt und später zentral (THM-CAS) angebunden
  werden. Die dafür nötigen Vorgaben bzw.\ das Datenschema wurden von der
  zuständigen Stelle nicht bereitgestellt, sodass die Umsetzung nicht
  weiterverfolgt wurde; die Anwendung bleibt ohne Rechteeinschränkung.
  \item \textbf{KI-Funktion.} Eine im Lastenheft angedachte KI-Unterstützung
  wurde im Projektverlauf --- mit Blick auf Umfang und die verkleinerte
  Teamgröße --- nicht implementiert.
\end{itemize}

\subsection{Bewusst zurückgestellt}
\begin{itemize}
  \item \textbf{Modifikatoren.} Die Entität und Verknüpfung (\texttt{modifier},
  \texttt{item\_modifier}) sind im Datenmodell vorhanden, es gibt aber bewusst
  \emph{keine} UI. Auf Wunsch des Betreuers wurden die Modifikatoren zunächst
  ausgeklammert, da noch nicht abschließend geklärt ist, wie die automatische
  Variantenerzeugung fachlich ausgestaltet sein soll.
\end{itemize}

\subsection{Bekannte technische Grenzen}
\begin{itemize}
  \item \textbf{Ein Kollektions-Elternteil je Item.} Die zugrunde liegende
  Tabelle \texttt{item\_collection\_sub\_item} erlaubt theoretisch, dass ein Item
  in mehreren Kollektionen liegt. Die Baum-Oberfläche stellt ein Item jedoch an
  genau einer Stelle dar und setzt damit praktisch einen einzigen
  Kollektions-Elternteil voraus.
  \item \textbf{Umfang des zentralen Stores.} Der Pinia-Store ist mit der Zeit
  stark gewachsen. Eine Aufteilung in fachliche Composables (Baum/Kollektionen,
  Inhalte, Tags, Suche) würde die Wartbarkeit weiter verbessern.
\end{itemize}

\subsection{Ausblick}
Für kommende Iterationen bieten sich an: das vollständige Cluster-Deployment
(sobald eine Infrastruktur bereitsteht), die zentrale Anmeldung (THM-CAS) mit
darauf aufbauendem Rechtemodell, die Ausgestaltung der Modifikatoren, ein
weiterer Ausbau der automatisierten Testabdeckung sowie die Modularisierung des
Stores.

% ══════════════════════════════════════════════════════════════════════════════
\section{Fazit}

Das Projekt setzt das vorgegebene Datenmodell in einer vollständig lauffähigen
Drei-Schichten-Anwendung um. Die zentralen fachlichen Konzepte --- Metadaten,
Inhalte in mehreren Formaten, geordnete und ungeordnete Kollektionen, horizontale
Varianten, Validatoren, hierarchische Tags mit Suche sowie Darstellungs-Templates
--- sind implementiert und über eine benutzerfreundliche Oberfläche zugänglich.
Tragende Entwurfsentscheidungen, insbesondere die saubere Trennung von vertikalen
und horizontalen Beziehungen und die Entkopplung der Datenzugriffsschicht über
ein Adapter-Muster, sorgen für fachliche Klarheit und Erweiterbarkeit. Mit der
durchgehenden Typsicherheit, der regelbasierten Strukturvalidierung und der
Container-Grundlage besteht eine solide Basis, auf der --- sobald die externen
Voraussetzungen (Infrastruktur, Auth-Vorgaben) gegeben sind --- die verbleibenden
Punkte aufsetzen können.

% ══════════════════════════════════════════════════════════════════════════════
\begin{thebibliography}{9}
\bibitem{vue} Vue.js --- The Progressive JavaScript Framework. \url{https://vuejs.org}
\bibitem{vuetify} Vuetify --- Vue Component Framework. \url{https://vuetifyjs.com}
\bibitem{pinia} Pinia --- The Vue Store. \url{https://pinia.vuejs.org}
\bibitem{vuedraggable} Vue.Draggable. \url{https://github.com/SortableJS/vue.draggable.next}
\bibitem{codemirror} CodeMirror 6. \url{https://codemirror.net}
\bibitem{vite} Vite --- Next Generation Frontend Tooling. \url{https://vitejs.dev}
\bibitem{spring} Spring Boot Reference Documentation. \url{https://spring.io/projects/spring-boot}
\bibitem{postgres} PostgreSQL 16 Documentation. \url{https://www.postgresql.org/docs/16/}
\bibitem{k8s} Kubernetes Documentation. \url{https://kubernetes.io/docs/}
\bibitem{argocd} Argo CD Documentation. \url{https://argo-cd.readthedocs.io}
\bibitem{traefik} Traefik Proxy Documentation. \url{https://doc.traefik.io/traefik/}
\bibitem{thm} THM E-Learning-Plattform --- Onboarding-Leitfaden (internes Dokument).
\end{thebibliography}

% ══════════════════════════════════════════════════════════════════════════════
\newpage
\section*{Anhang}
\addcontentsline{toc}{section}{Anhang}

\subsection*{A. Lokale Ausführung}
\begin{verbatim}
docker compose up --build
# Frontend: http://localhost:8085   Backend: 8080 (intern)   DB: 5432
# Dummy-Modus (ohne Backend):  npm run dev:dummy
\end{verbatim}

\subsection*{B. Deployment-Konfiguration (Auszug values.yaml, Backend)}
\begin{verbatim}
name: pwi-aufgabendatenbank-backend
image:
  repository: ghcr.io/.../pwi-aufgabendatenbank-backend
  tag: latest
replicas: 1
containerPort: 8080
route:
  path: /api
\end{verbatim}

\subsection*{C. Screenshots der Anwendung}
\platzhalter{[Screenshot: Strukturbaum mit Kollektionen und Aufgaben]}
\platzhalter{[Screenshot: Editor mit Inhalten, Tags und Varianten]}
\platzhalter{[Screenshot: obere Filterleiste (Tag-Filter + Suche)]}

\subsection*{D. Weiteres}
\begin{itemize}
  \item ER-Diagramm des Datenmodells \textit{(hier einfügen)}.
  \item GitHub-Repository: \texttt{PWI-Aufgabendatenbank}.
\end{itemize}

% Hinweis: Screenshots einfuegen -> \platzhalter{...} ersetzen durch z. B.
% \begin{figure}[h]\centering\includegraphics[width=0.9\linewidth]{bild.png}
% \caption{...}\end{figure}

\end{document}
